\section{Conclusion}

The study of technological evolution has methods, theories, and abundant raw material, but it still lacks one important thing: a shared empirical substrate for versioned material-culture lineages. Patent networks provide part of that substrate, but only part. To understand why technologies change, it is not enough to record formal prior art. We also need structured representations of the broader selection environment in which technologies are judged, adopted, modified, and abandoned.

The Technological Phylogenetic Repository is proposed as that missing infrastructure. By combining commit-based lineage modeling, explicit pressure annotation, and multi-agent construction from heterogeneous evidence, it offers a path toward a more empirical science of technological evolution. The proposal is ambitious, but it is also testable. A focused pilot can determine whether the framework is historically credible, technically tractable, and analytically worthwhile.

If successful, TPRs could provide a common substrate for work spanning cultural evolution, computational memetics, innovation studies, and the history of technology. If unsuccessful, the field would still benefit from the attempt because it would clarify which parts of technological history can be formalized and which resist reduction. Either outcome would move the field forward.
